
\documentclass{article}
\usepackage{colortbl}
\usepackage{makecell}
\usepackage{multirow}
\usepackage{supertabular}

\begin{document}

\newcounter{utterance}

\centering \large Interaction Transcript for game `dond', experiment `semi', episode 11 with Qwen3.6{-}35B{-}A3B{-}FP8{-}without{-}reasoning.
\vspace{24pt}

{ \footnotesize  \setcounter{utterance}{1}
\setlength{\tabcolsep}{0pt}
\begin{supertabular}{c@{$\;$}|p{.15\linewidth}@{}p{.15\linewidth}p{.15\linewidth}p{.15\linewidth}p{.15\linewidth}p{.15\linewidth}}
    \# & $\;$A & \multicolumn{4}{c}{Game Master} & $\;\:$B\\
    \hline

    \theutterance \stepcounter{utterance}  
    & & \multicolumn{4}{p{0.6\linewidth}}{
        \cellcolor[rgb]{0.9,0.9,0.9}{
            \makecell[{{p{\linewidth}}}]{
                \texttt{\tiny{[P1$\langle$GM]}}
                \texttt{Igrate kooperativnu igru pregovaranja u kojoj se trebate dogovoriti s drugim igračem o tome kako podijeliti skup stavki između vas.} \\
\\ 
\texttt{Pravila:} \\
\texttt{(a) Vi i drugi igrač dobivate skup stavki. Svakom od vas također se daje tajna funkcija vrijednosti koja predstavlja koliko cijenite svaku vrstu predmeta.} \\
\texttt{(b) Razmjenjujete poruke s drugim igračem kako biste se dogovorili o tome tko dobiva koje stavke. Svaki može poslati maksimalno 5 poruka ili prekinuti igru ranije dajući tajni prijedlog u bilo kojem trenutku.} \\
\texttt{(c) Od vas se traži da podnese tajni prijedlog koji označava stavke koje želite, formatirane u uglatim zagradama kako slijedi: "{[}Prijedlog: <broj> <naziv stavke>, <broj> <naziv stavke>, <...>{]}"} \\
\texttt{(d) Ako su vaši prijedlozi komplementarni, tj. ima dovoljno stavki za ispunjavanje oba prijedloga, svaki igrač dobiva rezultat na temelju zbroja vrijednosti stavki koje je dobio. U suprotnom, obojica dobivate nula bodova.} \\
\texttt{(e) Vaš cilj je maksimizirati zbroj vašeg rezultata i rezultata drugog igrača.} \\
\\ 
\texttt{Počnimo.} \\
\\ 
\texttt{Skup dostupnih stavki je:} \\
\\ 
\texttt{1 knjiga, 1 stijena, 1 torba, 2 lonci} \\
\\ 
\texttt{Vaša tajna funkcija vrijednosti je:} \\
\\ 
\texttt{knjiga: 7, stijena: 2, torba: 1, lonac: 0} \\
\\ 
\texttt{VAŽNO: Vaše poruke, osim ako nisu tajni prijedlog, šalju se izravno drugom igraču, pa ne uključujte nikakav odgovor na pravila niti tekst koji najavljuje vašu poruku. Za davanje tajnog prijedloga koristite navedeni format. Ne koristite uglate zagrade pri komuniciranju s drugim igračem ili će se tumačiti kao vaš tajni prijedlog.} \\
            }
        }
    }
    & \\ \\

    \theutterance \stepcounter{utterance}  
    & \multicolumn{4}{p{0.6\linewidth}}{
        \cellcolor[rgb]{0.9,0.9,0.9}{
            \makecell[{{p{\linewidth}}}]{
                \texttt{\tiny{[P1$\rangle$GM]}}
                \texttt{Pozdrav! Vidim da imamo skup od 1 knjige, 1 stijene, 1 torbe i 2 lonca.} \\
\\ 
\texttt{Kako bismo maksimizirali ukupnu vrijednost, važno je da stavke koje visoko cijenimo završe kod onoga tko ih više cijeni. Ja sam naveo da mi knjiga vrijedi 7, a ostale stavke vrlo malo ili nula. Pretpostavljam da i vi imate slične preferencije ili barem da knjiga ima visoku vrijednost za nekoga od nas.} \\
\\ 
\texttt{Predlažem da knjiga pripadne onome tko je više cijeni. Ako vi knjigu cijenite više od mene (ili jednako), predlažem da vi dobijete knjigu. Stijena i torba imaju niske vrijednosti, a lonci su za mene bezvrijedni (0), pa ih možemo podijeliti kako god vama više odgovara, sve dok ne oduzmete knjigu ako je vi ne želite.} \\
\\ 
\texttt{Moja prva poruka:} \\
\texttt{"Predlažem da knjiga pripadne onome tko je više cijeni. Ako vi knjigu cijenite više, neka ona bude vaša. Stijenu i torbu možemo podijeliti po želji, a lonci su za mene bezvrijedni, pa ih možete zadržati ili dati meni, ovisno o vašoj procjeni. Što kažete?"} \\
            }
        }
    }
    & & \\ \\

    \theutterance \stepcounter{utterance}  
    & & & \multicolumn{2}{p{0.3\linewidth}}{
        \cellcolor[rgb]{0.9,0.9,0.9}{
            \makecell[{{p{\linewidth}}}]{
                \texttt{\tiny{[GM$|$GM]}}
                \texttt{continue} \\
            }
        }
    }
    & & \\ \\

    \theutterance \stepcounter{utterance}  
    & & & \multicolumn{2}{p{0.3\linewidth}}{
        \cellcolor[rgb]{0.9,0.9,0.9}{
            \makecell[{{p{\linewidth}}}]{
                \texttt{\tiny{[GM$|$GM]}}
                \texttt{Pozdrav! Vidim da imamo skup od 1 knjige, 1 stijene, 1 torbe i 2 lonca.} \\
\\ 
\texttt{Kako bismo maksimizirali ukupnu vrijednost, važno je da stavke koje visoko cijenimo završe kod onoga tko ih više cijeni. Ja sam naveo da mi knjiga vrijedi 7, a ostale stavke vrlo malo ili nula. Pretpostavljam da i vi imate slične preferencije ili barem da knjiga ima visoku vrijednost za nekoga od nas.} \\
\\ 
\texttt{Predlažem da knjiga pripadne onome tko je više cijeni. Ako vi knjigu cijenite više od mene (ili jednako), predlažem da vi dobijete knjigu. Stijena i torba imaju niske vrijednosti, a lonci su za mene bezvrijedni (0), pa ih možemo podijeliti kako god vama više odgovara, sve dok ne oduzmete knjigu ako je vi ne želite.} \\
\\ 
\texttt{Moja prva poruka:} \\
\texttt{"Predlažem da knjiga pripadne onome tko je više cijeni. Ako vi knjigu cijenite više, neka ona bude vaša. Stijenu i torbu možemo podijeliti po želji, a lonci su za mene bezvrijedni, pa ih možete zadržati ili dati meni, ovisno o vašoj procjeni. Što kažete?"} \\
            }
        }
    }
    & & \\ \\

    \theutterance \stepcounter{utterance}  
    & & \multicolumn{4}{p{0.6\linewidth}}{
        \cellcolor[rgb]{0.9,0.9,0.9}{
            \makecell[{{p{\linewidth}}}]{
                \texttt{\tiny{[GM$\rangle$P2]}}
                \texttt{Igrate kooperativnu igru pregovaranja u kojoj se trebate dogovoriti s drugim igračem o tome kako podijeliti skup stavki između vas.} \\
\\ 
\texttt{Pravila:} \\
\texttt{(a) Vi i drugi igrač dobivate skup stavki. Svakom od vas također se daje tajna funkcija vrijednosti koja predstavlja koliko cijenite svaku vrstu predmeta.} \\
\texttt{(b) Razmjenjujete poruke s drugim igračem kako biste se dogovorili o tome tko dobiva koje stavke. Svaki može poslati maksimalno 5 poruka ili prekinuti igru ranije dajući tajni prijedlog u bilo kojem trenutku.} \\
\texttt{(c) Od vas se traži da podnese tajni prijedlog koji označava stavke koje želite, formatirane u uglatim zagradama kako slijedi: "{[}Prijedlog: <broj> <naziv stavke>, <broj> <naziv stavke>, <...>{]}"} \\
\texttt{(d) Ako su vaši prijedlozi komplementarni, tj. ima dovoljno stavki za ispunjavanje oba prijedloga, svaki igrač dobiva rezultat na temelju zbroja vrijednosti stavki koje je dobio. U suprotnom, obojica dobivate nula bodova.} \\
\texttt{(e) Vaš cilj je maksimizirati zbroj vašeg rezultata i rezultata drugog igrača.} \\
\\ 
\texttt{Počnimo.} \\
\\ 
\texttt{Skup dostupnih stavki je:} \\
\\ 
\texttt{1 knjiga, 1 stijena, 1 torba, 2 lonci} \\
\\ 
\texttt{Vaša tajna funkcija vrijednosti je:} \\
\\ 
\texttt{knjiga: 0, stijena: 1, torba: 7, lonac: 1} \\
\\ 
\texttt{VAŽNO: Vaše poruke, osim ako nisu tajni prijedlog, šalju se izravno drugom igraču, pa ne uključujte nikakav odgovor na pravila niti tekst koji najavljuje vašu poruku. Za davanje tajnog prijedloga koristite navedeni format. Ne koristite uglate zagrade pri komuniciranju s drugim igračem ili će se tumačiti kao vaš tajni prijedlog.} \\
\\ 
\texttt{Pozdrav! Vidim da imamo skup od 1 knjige, 1 stijene, 1 torbe i 2 lonca.} \\
\\ 
\texttt{Kako bismo maksimizirali ukupnu vrijednost, važno je da stavke koje visoko cijenimo završe kod onoga tko ih više cijeni. Ja sam naveo da mi knjiga vrijedi 7, a ostale stavke vrlo malo ili nula. Pretpostavljam da i vi imate slične preferencije ili barem da knjiga ima visoku vrijednost za nekoga od nas.} \\
\\ 
\texttt{Predlažem da knjiga pripadne onome tko je više cijeni. Ako vi knjigu cijenite više od mene (ili jednako), predlažem da vi dobijete knjigu. Stijena i torba imaju niske vrijednosti, a lonci su za mene bezvrijedni (0), pa ih možemo podijeliti kako god vama više odgovara, sve dok ne oduzmete knjigu ako je vi ne želite.} \\
\\ 
\texttt{Moja prva poruka:} \\
\texttt{"Predlažem da knjiga pripadne onome tko je više cijeni. Ako vi knjigu cijenite više, neka ona bude vaša. Stijenu i torbu možemo podijeliti po želji, a lonci su za mene bezvrijedni, pa ih možete zadržati ili dati meni, ovisno o vašoj procjeni. Što kažete?"} \\
            }
        }
    }
    & \\ \\

    \theutterance \stepcounter{utterance}  
    & & & \multicolumn{4}{p{0.6\linewidth}}{
        \cellcolor[rgb]{0.9,0.9,0.9}{
            \makecell[{{p{\linewidth}}}]{
                \texttt{\tiny{[GM$\langle$P2]}}
                \texttt{"Razumijem vašu poruku, ali moram vas ispraviti u vezi s mojim vrijednostima. U pravilima je navedeno da je moja tajna funkcija vrijednosti: knjiga: 0, stijena: 1, torba: 7, lonac: 1.} \\
\\ 
\texttt{To znači da mi knjiga ne vrijedi ništa, dok mi torba vrijedi 7 bodova. Stijena i lonci vrijede po 1 bod.} \\
\\ 
\texttt{Budući da je torba za mene najvrijednija stavka (7 bodova), a za vas vjerojatno ima neku vrijednost (iako ne znam vašu točnu funkciju), predlažem sljedeće kako bismo maksimizirali ukupnu vrijednost:} \\
\\ 
\texttt{1. Ja dobivam torbu (vrijednost 7 za mene).} \\
\texttt{2. Vi dobivate knjigu (vrijednost 0 za mene, ali možda više za vas).} \\
\texttt{3. Stijenu i dva lonca možemo podijeliti tako da svatko dobije ono što više cijeni.} \\
\\ 
\texttt{Budući da ne znam vaše vrijednosti, predlažem da vi zadržite knjigu, a ja uzmem torbu. Za ostale stavke (stijena, 2 lonca), predlažem da ih podijelimo jednako ili da ih dobijete vi ako ih više cijenite.} \\
\\ 
\texttt{Moj tajni prijedlog za sada je:} \\
\texttt{{[}Prijedlog: 1 torba{]}"} \\
            }
        }
    }
     \\ \\

    \theutterance \stepcounter{utterance}  
    & & & \multicolumn{2}{p{0.3\linewidth}}{
        \cellcolor[rgb]{0.9,0.9,0.9}{
            \makecell[{{p{\linewidth}}}]{
                \texttt{\tiny{[GM$|$GM]}}
                \texttt{proposal} \\
            }
        }
    }
    & & \\ \\

    \theutterance \stepcounter{utterance}  
    & & & \multicolumn{2}{p{0.3\linewidth}}{
        \cellcolor[rgb]{0.9,0.9,0.9}{
            \makecell[{{p{\linewidth}}}]{
                \texttt{\tiny{[GM$|$GM]}}
                [0, 0, 1, 0]
            }
        }
    }
    & & \\ \\

    \theutterance \stepcounter{utterance}  
    & & \multicolumn{4}{p{0.6\linewidth}}{
        \cellcolor[rgb]{0.9,0.9,0.9}{
            \makecell[{{p{\linewidth}}}]{
                \texttt{\tiny{[P1$\langle$GM]}}
                \texttt{Drugi igrač je poslao tajni prijedlog. Sada je na tebi red da pošalješ vlastiti prijedlog.} \\
            }
        }
    }
    & \\ \\

    \theutterance \stepcounter{utterance}  
    & \multicolumn{4}{p{0.6\linewidth}}{
        \cellcolor[rgb]{0.9,0.9,0.9}{
            \makecell[{{p{\linewidth}}}]{
                \texttt{\tiny{[P1$\rangle$GM]}}
                \texttt{{[}Prijedlog: 1 knjiga, 1 torba, 2 lonca{]}} \\
            }
        }
    }
    & & \\ \\

    \theutterance \stepcounter{utterance}  
    & & & \multicolumn{2}{p{0.3\linewidth}}{
        \cellcolor[rgb]{0.9,0.9,0.9}{
            \makecell[{{p{\linewidth}}}]{
                \texttt{\tiny{[GM$|$GM]}}
                \texttt{proposal must include only valid item types} \\
            }
        }
    }
    & & \\ \\

\end{supertabular}
}

\end{document}
